\newif\ifglobal

%%%%%%%%%%%%%%%%%%%%%%%
%%% Compile to view %%%
%%%%%%%%%%%%%%%%%%%%%%%

%%%%%%%%%%%%%%%%%%%%%%%%%%%%%%%%%%%%%%%%%%%%%%%%%%%
%%% Readme:                                     %%%
%%% https://www.overleaf.com/read/bwdjvfjksvjy/ %%%
%%%%%%%%%%%%%%%%%%%%%%%%%%%%%%%%%%%%%%%%%%%%%%%%%%%

% >>> M?: marks a module setting number or important notice

% >>> M4: Set {./relative-path} to external files for this module <<<
\def \modulefiles {./34-RNG}
% To add an external file, use:
% (Replace '---' with actual filename)
% '\input{\modulefiles/tables/---.tex}' for tables
% '\includegraphics{\modulefiles/figures/---}' for figures
% '\subfile{\modulefiles/---}' for register files


% >>> M1: This module is in global project? <<<
% 'YES' - uncomment; 'NO' - comment out
\globaltrue

\ifglobal
    \documentclass[main\_\_CN.tex]{subfiles}

\else
    % Font "FandolSong-Regular" used by ctex does not have GSUB table.
    % It causes warnings that do not affect compilation and can be ignored.
    \PassOptionsToPackage{quiet}{fontspec}

    \documentclass[openany, 10pt]{book}

    % >>> M2: In ./chip-spec-settings/preamble-trm-module, also find
    %         settings for visibility of todo notes, watermark <<<
    \usepackage{preamble-trm-module}


    % Enable Chinese typesetting
    \usepackage{ctex}

    % >>> M3: Temporary labels in Chinese
    %         you can add your own labels there <<<
    \myexternaldocument{temp-labels-cn}

\fi %\ifglobal


\setcounter{secnumdepth}{4}



% >>> M5: If this module needs any updates in
% './chip-spec-settings/preamble-trm-module.sty'
% (include additional package, etc.), contact your module PM
% or create the file './chip-spec-settings/changelog.tex' make the
% necessary updates yourself and document them in this file <<<

\raggedright


\begin{document}


% Sets language into which pre-defined bits of text are translated
\selectlanguage{chinese}

% Do not delete this block
\ifglobal
% Add the following front matter if
% this module is outside of global project
\else
    \InputIfFileExists{readme}{}{}
    {\let\clearpage\relax \listoftodos}
    {\let\clearpage\relax \listofdonetodos}
    \tableofcontents
    %\listoftables
    %\listoffigures
\fi


%%%%%%%%%%%%%%%%%%%%%%%%%
%%% TRM Module Begins %%%
%%%%%%%%%%%%%%%%%%%%%%%%%

\hypertarget{rng}{}
\chapter{随机数发生器 (RNG)}\label{mod:rng}

\section{概述}

\chipname{} 内置一个真随机数发生器，其生成的 32 位随机数可作为加密等操作的基础。

\section{主要特性}
\label{sec:rng-features}

\chipname{} 的随机数发生器可通过物理过程而非算法生成真随机数，所有生成的随机数在特定范围内出现的概率完全一样。

\section{功能描述}\label{sec:rng-func-descr}

系统可以从随机数发生器的寄存器 \hyperref[regdesc:LPSYSREGRNGDATAREG]{LPSYSREG\_RNG\_DATA\_REG} 中读取随机数，每个读到的 32 位随机数都是真随机数，噪声源为系统中的\textbf{热噪声}和\textbf{异步时钟}。

\begin{itemize}
    \item \textbf{热噪声}可以来自 SAR ADC。当芯片的 SAR ADC 工作时，就会产生比特流，并通过异或 (XOR) 逻辑运算作为随机数种子进入随机数生成器。
    \item RC\_FAST\_CLK 是一种\textbf{异步时钟源}，会产生电路亚稳态。这种亚稳态也可以作为随机数种子，进入随机数生成器。关于 RC\_FAST\_CLK，详见以下 \hyperref[others:rng-note-clock-seed]{说明}。
\end{itemize}
\vspace{2em}
\begin{figure}[H]
    \centering
    \includegraphics[width=1.0\textwidth]{\modulefiles/figures/34-RNG-noise-source}
    \caption{噪声源}
    \label{fig:rng-noise-source}
\end{figure}
当 SAR ADC 打开时，每个 RC\_FAST\_CLK 时钟周期内（来自内部 RC 振荡器，详见 \ref{mod:resclk} \textit{\nameref{mod:resclk}} 章节），随机数发生器将获得 1 位的熵。因此，为了获得最大的熵值，建议读取 \hyperref[regdesc:LPSYSREGRNGDATAREG]{LPSYSREG\_RNG\_DATA\_REG} 寄存器时的速率不超过 1 MHz。


\section{编程指南}

在使用 \chipname{} 的随机数生成器时，应该至少保证 SAR ADC 或 RC\_FAST\_CLK 处于使能状态，否则可能会导致产生伪随机数，应注意避免。其中，
\begin{itemize}
    \item 关于 SAR ADC，详见 \ref{mod:adc} \textit{\nameref{mod:adc}} 章节。
    \item RC\_FAST\_CLK 可通过设置 \hyperref[regdesc:PMUHPSLEEPLPCKPOWERREG]{PMU\_HP\_SLEEP\_LP\_CK\_POWER\_REG} 寄存器中的 \hyperref[fielddesc:PMU_HPSLEEPXPDFOSCCLK]{PMU\_HP\_SLEEP\_XPD\_FOSC\_CLK} 位使能。关于 RC\_FAST\_CLK，详见以下 \hyperref[others:rng-note-clock-seed]{说明}。
    \item 通过设置 \hyperref[regdesc:LPPERICLKENREG]{LPPERI\_CLK\_EN\_REG} 寄存器中的 \hyperref[fielddesc:LPPERICKENRNG]{LPPERI\_CK\_EN\_RNG} 位，对 RC\_FAST\_CLK 进行门控。
\end{itemize}
\begin{tiplisting}
    \label{others:rng-note-clock-seed}
    RC\_FAST\_CLK 时钟仅可以提高随机数发生器的墒值。然而，为了保证随机数发生器可以获得足够大的熵值，仍建议在使用随机数发生器时至少保证 SAR ADC 处于工作状态。
\end{tiplisting}

在使用随机数生成器时，请多次读取 \hyperref[regdesc:LPSYSREGRNGDATAREG]{LPSYSREG\_RNG\_DATA\_REG} 寄存器的值，直至获得足够多的随机数。在读取寄存器时，注意控制速率不要超过上方第 \ref{sec:rng-func-descr} 小节的介绍。


% >>> If your module has no registers, delete sample text
%       for Sections Register Summary and Registers below <<<

%%%%%%%%%%%%%%%%%%%%%%%%
%%% Register Summary %%%
%%%%%%%%%%%%%%%%%%%%%%%%

% >>> Update the sample text below and insert
%     the info on register summary into the subfile <<<

\newpage

\hypertarget{rng-reg-summ}{}
\section{寄存器列表}
\label{sec:rng-reg-summ}


请查看章节 \hyperref[glossary-access-types]{\textit{寄存器的访问类型}}，了解“访问”列缩写的含义。

\begin{table}[H]
    \centering
    \begin{tabular}{ | l | l | l | l |  }
    \hline
    \rowcolor{lightgray}
    \textbf{名称} & \textbf{描述} & \textbf{地址} & \textbf{访问} \\
    \hline
    \hyperref[regdesc:LPSYSREGRNGDATAREG]{LPSYSREG\_RNG\_DATA\_REG} & 随机数源数据 & 0x5011\_01A4 & RO \\ \hline
    \end{tabular}
\end{table}



%%%%%%%%%%%%%%%%%
%%% Registers %%%
%%%%%%%%%%%%%%%%%

% >>> Update the sample text below and insert
%      the info on registers into the subfile <<<

\newpage
\section{寄存器}


关于保留 (reserved) 域的处理，请查看章节 \hyperref[programming-reserved-field]{如何配置寄存器的保留域}。

\begin{register}{H}{LPSYSREG\_RNG\_DATA\_REG}{0x5011\_01A4}\label{regdesc:LPSYSREGRNGDATAREG}
\regfield{LPSYSREG\_RNG\_DATA}{32}{0}{{0x00000000}}%
\reglabel{Reset}\regnewline%
\begin{regdesc}\begin{reglist}
\label{fielddesc:LPSYSREGRNGDATA}\item [LPSYSREG\_RNG\_DATA] 随机数源数据。(RO)
\end{reglist}\end{regdesc}
\end{register}
\end{document}
